% !PRINTED 382
\weberpage{382}

%% [running head, not body text -- commented out 2026-08-07] Page furniture (running
%% head, folio, signature) is excluded from this edition by standing policy.
%% \noindent 382 \hfill \textsc{Eleventh Section.} \hfill \S.\,123.

The reduced numbers whose images lie on the boundary of the fundamental
region are equivalent in pairs, namely

\noindent a)\quad $\omega = -\tfrac{1}{2} + i\eta$ \quad and \quad $\omega_1 = \omega + 1 = \tfrac{1}{2} + i\eta$

\noindent ($a$ and $a'$ in the figure),

\noindent b)\quad $\omega = -\xi + i\eta$ \quad and \quad $\omega_1 = -\dfrac{1}{\omega} = \xi + i\eta$

\noindent when $\xi^2+\eta^2 = 1$ ($b$ and $b'$ in the figure). But there holds further the
theorem:

\noindent 2.\quad \gesperrt{Apart from the cases a), b), no two reduced numbers are equivalent.}

For let us assume two reduced numbers that are not identical,

\weq{11}{\omega = \xi + \eta i, \quad \omega_1 = \xi_1 + \eta_1 i,}

and suppose, which does not impair the generality,

\weq{12}{\eta_1 \geqq \eta,}

then it follows from formula (4):

\weq{13}{\eta_1 = \dfrac{\eta}{(\gamma\xi+\delta)^2 + \gamma^2\eta^2},}

hence, on account of (12):

\weq{14}{(\gamma\xi+\delta)^2 + \gamma^2\eta^2 \leqq 1,}

hence also, by (3):

\weq{15}{1 \geqq \gamma^2\eta^2 \geqq \tfrac{3}{4}\gamma^2;}

thus there are two possibilities:
\[
\alpha)\ \gamma = 0, \qquad \beta)\ \gamma = \pm 1.
\]

In the case $\alpha$) we have $\alpha\delta = 1$, hence $\alpha = \delta = \pm 1$; from (12) and (13)
it results that $\eta_1 = \eta$, and further from (4):
\[
\xi_1 = \xi \pm \beta.
\]
But this is compatible with the conditions (2) only when
either $\xi_1 = \xi$, $\beta = 0$ and hence $\omega_1$ is identical with $\omega$, or when
$\beta = \pm 1$ and $\xi_1 = -\xi = \pm \tfrac{1}{2}$, that is, in the case a).

In the case $\beta$) we have by (14) $(\xi \pm \delta)^2 \leqq 1$, and hence \gesperrt{either}:
$\delta = 0$, $\beta\gamma = -1$ and $\xi^2+\eta^2 = 1$, $\eta_1 = \eta$, further by (4)
$\xi_1 = \pm\alpha - \xi$; hence, since $\xi_1 = \xi$ is excluded, $\xi_1 = -\xi$,
$\alpha = 0$, which is the exceptional case b),

\gesperrt{or}: $\delta = \pm 1$, hence $(\xi\pm 1)^2 + \eta^2 \leqq 1$, and consequently by (3)
$(\xi\pm 1)^2 \leqq \tfrac{1}{4}$. But since $(\xi\pm 1)$ must by (2) be at least
equal to $\tfrac{1}{2}$ in absolute value, it follows from this that $\xi = \mp\tfrac{1}{2}$,
$\eta^2 = \tfrac{3}{4}$; hence by (13) $\eta_1^2 = \eta^2 = \tfrac{3}{4}$, and by
%CONT>>
